\documentclass[9pt, aspectratio=1610]{beamer}
%\documentclass{article}
%\usepackage{beamerarticle}
\mode<article>
{
	\usepackage[a4paper]{geometry}
}



\edef\tlCourse{Промышленное программирование}
\edef\tlPartName{Модуль 1. Разработка модульных программ}
\edef\tlPartNumber{1}
\edef\tlThemeName{Разработка модульных программ на С++}
\edef\tlThemeNumberLocal{2}
\edef\tlThemeNumberGlobal{2}


%\documentclass[10pt, handout]{beamer}
%\setbeameroption{show notes}

%\documentclass[10pt, a4paper]{article}
%\usepackage{beamerarticle}



%Для XeLatex/+
%\usepackage{polyglossia}
%\setdefaultlanguage{russian}
%\setotherlanguage{english}
%%\setkeys{russian}{babelshorthands=true} 

\usepackage[russian]{babel}



\usepackage{fontspec}

\setmainfont{Times New Roman} 
\setsansfont{Arial} 
\setmonofont{Consolas}  



\mode<article>{
	
	\usepackage{hyperref}
	
}
\mode<presentation>{
	
	%\usetheme{Antibes}
	\usetheme{Berlin}
	
	\usefonttheme{professionalfonts} 
	\usefonttheme{serif} % default family is serif
	
	\usepackage{fontspec}
	\setmonofont{Cascadia Code SemiLight} %CMU Typewriter Text
	\setromanfont{DejaVu Serif}
	\setsansfont{DejaVu Sans}
	%\usecolortheme{spruce} %зеленая, плохой цвет в заголовках 
	%\usecolortheme{albatross} %синяя, пхоло виден черный цвет
	
	\setbeamercovered{dynamic}
	
}

\newlength{\parskipbackup}
\setlength{\parskipbackup}{\parskip}
\newlength{\parindentbackup}
\setlength{\parindentbackup}{\parindent}

\let\notebackup\note
\renewcommand{\note}[1]{\notebackup{%
		\setlength{\parindent}{3ex}%
		\setlength{\parskip}{0.5ex}%
		{\small{}#1}%
		\setlength{\parskip}{\parskipbackup}%
		\setlength{\parindent}{\parindentbackup}%
	}%
}

\mode<presentation>%
{
	
	%\setbeameroption{second mode text on second screen=right}
	%\setbeameroption{show notes on second screen}
	%\setbeameroption{show notes}
	\setbeamertemplate{note page}{
		
		%\vspace{1em}\insertframetitle
		
		%		\begin{wrapfigure}[12]{l}[0cm]{0.5\textwidth}
			%			\fbox{ \resizebox{0.45\textwidth}{!}{\insertslideintonotes{1}} }
			%		\end{wrapfigure}		
		\insertnote}
	
	
}


\newcommand{\MP}[1]{\mode<presentation>{#1} }
\newcommand{\MA}[1]{\mode<article>{#1} }

\newcommand{\ABS}[1]{\left| #1 \right|}
%\newcommand{\ABS}[1]{\mid #1 \mid}

\newcommand{\HREF}[2]{{\color{blue}\underline{\href{#1}{#2}}}}

\setbeamertemplate{caption}[numbered]


\usepackage{array}
\usepackage{tabularx}

\usepackage{verbatim}

\usepackage{ifthen}



\usepackage{tikz}
\usetikzlibrary{arrows.meta}
\usetikzlibrary{calc}
\usetikzlibrary{decorations}
\usetikzlibrary{decorations.pathreplacing}
\usetikzlibrary{matrix}
\usetikzlibrary{fit}

\newcommand{\rememb}[1]{\parbox{0pt}{}\tikz[remember picture,baseline=-0.5ex]{\draw node[inner sep=0pt, outer sep=0pt] (#1){\strut};}}



\usepackage{tikz-3dplot}
\usepackage{environ}
\usepackage{animate}





\usepackage{xcolor}


\usepackage{amsmath} %математические формулы
%\usepackage{amsfonts} %математические шрифты


\usepackage[e]{esvect}  %Красивая стрелочка вектора
%\let\oldvv\vv
\newcommand{\VV}[1]{\vv{#1\mathstrut}}

\usepackage{graphicx} %работа с каритнками




\usepackage{array}
\usepackage{multirow}



\usepackage{caption}
\DeclareCaptionLabelSeparator{dot}{~---~}            %Разделитель номер рисунка
%\captionsetup[figure]{justification=centering,labelsep=dot, format=plain}                        %Подпись рис. центр
%\captionsetup[table]{justification=raggedleft,labelsep=dot, format=plain, singlelinecheck=false} %Подпись табл. слева
\captionsetup[lstlisting]{justification=raggedleft,labelsep=dot, format=plain, font=normal, singlelinecheck=false}                     %Подпись рис. центр

\usepackage{indentfirst} %отступ первой строки


\usepackage[svgnames]{xcolor}


\usepackage{hyperref}



\def\sectionname{Раздел}
\def\subsectionname{Подраздел}


\NewDocumentCommand\TC{m+m+m}{%	
	\begin{columns}
		\begin{column}{#1\textwidth}
			#2
		\end{column}
		\begin{column}{\fpeval{1-#1}\textwidth}
			#3
		\end{column}
	\end{columns}
}

\NewDocumentCommand\TCT{m+m+m}{%	
	\begin{columns}[T]
		\begin{column}{#1\textwidth}
			#2
		\end{column}
		\begin{column}{\fpeval{1-#1}\textwidth}
			#3
		\end{column}
	\end{columns}
}


\usepackage{newfile}

\edef\LectionNumber{0}
\edef\LectionTheme{0}

\let\oldsection\section
\let\oldsubsection\subsection


\AtBeginDocument
{
	\newoutputstream{CONTENT}
	\openoutputfile{\jobname.gvr}{CONTENT}
	
	%\expandafter\addtostream{CONTENT}{\noindent\textbf{\Large Лекция \LectionNumber~---~\LectionTheme}\unexpanded{\setcounter{SEC}{0}}\par}
	\addtostream{CONTENT}{\protect\noindent%
						  	\protect\setcounter{SEC}{0}%
							{%
							\protect\large% 
							\protect\bfseries{}%
							Лекция \LectionNumber~---~\LectionTheme% 
							}% 
							\protect\par}
}

\renewcommand{\section}[2][]{
	\ifstrempty{#1}%
	{\oldsection{#2}}%
	{\oldsection[#1]{#2}}	
	\addtostream{CONTENT}{%
			\protect\noindent%
			\protect\stepcounter{SEC}%
			\protect\setcounter{SUB}{0}%
			\protect\hspace{3ex}%
			\protect\theSEC ~ #1 \protect\par%
			}
}

\renewcommand{\subsection}[1]{
	\oldsubsection{#1}
	\addtostream{CONTENT}{%
			\protect\noindent%
			\protect\stepcounter{SUB}%
			\protect\hspace{6ex}%
			\protect\theSEC{}.\protect\theSUB ~ #1 \protect\par%
			}
}


\newcommand{\Strut}{{\Large\strut}}

\newcommand\scb[1]{\left( #1 \right)}

\newcommand{\LINK}[2]{%
	\qrcode[height=1cm]{#1}\  \HREF{#1}{\parbox{0.8\textwidth}{#2}} \\[0.5em]
}

\NewDocumentCommand{\lecdni}{}{\LectionNumber}
\author{Гаврилов Андрей Геннадьевич}
\newcommand{\regals}{кандидат технических наук, доцент}
\institute{Кафедра Информационных технологий и вычислительных систем \\МГТУ~<<СТАНКИН>>}
\lecture{История компьютерной графики}{kghistory}\subtitle{Компьютерная графика}


\makeatletter
\newcommand*{\overlaynumber}{\number\beamer@slideinframe}
\makeatother





\usepackage{cprotect}

\usepackage{qrcode}

\newcommand{\QRFRAME}{%
    \begin{frame}[plain, noframenumbering]    	
	
	\centering
	Трансляция презентации (во время очных лекций)    
	
	~
	
	{\Large \ttfamily  https://clck.ru/3D3Efj  }
	
	~
	
	\tikz\node[inner sep=0pt,rounded corners=5mm, clip]{\qrcode[height=0.4\textwidth]{https://clck.ru/3D3Efj}}; 
	
	~	
	{\small
	При просмотре презентации в PDF для отображения анимаций на слайдах необходимо использовать Acrobat Reader, KDE Okular, PDF-XChange, Foxit Reader, браузер Firefox. Для браузеров на движке Chrome (Edge, Яндекс, Opera,~\dots) необходимо использовать \HREF{https://chromewebstore.google.com/detail/pdf-viewer/oemmndcbldboiebfnladdacbdfmadadm?hl=ru&utm_source=ext_sidebar}{PDF.js} c опцией <<Enable active content (JavaScript) in PDFs>>. }
	
	\end{frame}%
}

\newcommand{\IG}[2][1]{\includegraphics[width=#1\textwidth]{#2}}


%===========Настройка листингов==================

\usepackage{fancyvrb}
\usepackage{minted}

\AtBeginEnvironment{minted}{%
	\renewcommand{\fcolorbox}[4][]{#4}}

\mode<presentation>
{
	\newminted[cppcode]{c++}{linenos, autogobble=true, style=vs,tabsize=2,fontsize=\small,escapeinside=\@\@,xleftmargin=0em, numbersep=0.5em}

	
	\newminted[cppresumecode]{c++}{linenos, firstnumber=last,  autogobble=true, style=vs,tabsize=2,fontsize=\small,escapeinside=\@\@,xleftmargin=0em, numbersep=0.5em}
	
	\newminted[cmakecode]{cmake}{linenos, autogobble=true, style=default,tabsize=2,fontsize=\footnotesize,escapeinside=\@\@,xleftmargin=0em, numbersep=0.5em}
	
	\newminted[cmdcode]{Batch}{linenos, autogobble=true, bgcolor=red,  style=lightbulb,tabsize=2,fontsize=\small,escapeinside=\@\@,xleftmargin=0em, numbersep=0.5em}
}

\mode<!presentation>
{
	\newminted[cppcode]{c++}{linenos,  frame=lines, style=vs,tabsize=3,fontsize=\small,escapeinside=\@\@,xleftmargin=2em}
	\newminted[cppresumecode]{c++}{linenos, firstnumber=last, frame=lines, style=vs,tabsize=3,fontsize=\small,escapeinside=\@\@,xleftmargin=2em}
	
	\newminted[cmakecode]{cmake}{linenos, autogobble=true, style=vs,tabsize=2,fontsize=\small,escapeinside=\@\@,xleftmargin=2em}
	
	\newminted[cmdcode]{batch}{linenos, autogobble=true, style=lightbulb,tabsize=2,fontsize=\small,escapeinside=\@\@,xleftmargin=0em, numbersep=0.5em}
}

\newenvironment{codeblock}[2][]
{
	\VerbatimEnvironment
	\begin{block}{#1\hfill{}#2}
		\begin{cppcode}% 
}
{		
	\end{cppcode}
	\end{block}
}

\newenvironment{codeblockresume}[2][]
{
	\VerbatimEnvironment
	\begin{block}{#1\hfill{}#2}
		\begin{cppresumecode}% 
}
{		
		\end{cppresumecode}
		\end{block}
}


%===================================================




\graphicspath{{Images/}{Images/\jobname/}}

\date{\today}

\renewcommand{\LectionNumber}{\tlThemeNumberGlobal}
\renewcommand{\LectionTheme}{Тема \tlThemeNumberGlobal \\ \tlThemeName}
\title{\LectionTheme}
\subtitle{\tlCourse}


\usepackage{fancyvrb}
\usepackage{minted}


\usetikzlibrary {shapes.geometric}
\usetikzlibrary{shapes.misc}
\usetikzlibrary{positioning}
\usetikzlibrary{calc}
\usetikzlibrary{trees}

\tikzset{	bblock/.style 	= {draw, minimum width=8em, minimum height=2em, text height=1.5ex,text depth=0.25ex},
	io/.style 		= {bblock, trapezium, trapezium left angle=70, trapezium right angle=110},
	op/.style 		= {bblock, rectangle}, 
	if/.style		= {bblock, diamond},
	terminator/.style		= {bblock, rounded rectangle},
	portal/.style 	={bblock,circle,minimum size=3em},
}

% Source - https://tex.stackexchange.com/a/23094
% Posted by egreg
% Retrieved 2026-02-16, License - CC BY-SA 3.0

\newcommand*{\fvtextcolor}[2]{\textcolor{#1}{#2}}
\newcommand*{\tc}[2]{\textcolor{#1}{#2}}


\begin{document}
	 
	\mode<presentation>
	{
		\input{template.tex}	
	}
	
	





\QRFRAME
	

	
\mode<presentation>
{	
	
		\frame{\maketitle}
		
		\begin{frame}\frametitle{План лекции}			
			\tableofcontents
		\end{frame}

	
}



\mode<article>
{
	\begin{center}
		
		\Large
		
		\title
		
		\LectionTheme
	\end{center}
}
	
\begin{frame}[fragile]\frametitle{Напишем математическую программу}
	
	\begin{block}{main.cpp}
		
		\begin{cppcode} 
			//include iostream cmath tuple exception clocale print
			
			std::tuple<double,double> quadratic_equation(double a, double b, double c)
			{
				double D = b*b-4*a*c;
				if (D>0)
				return std::make_tuple(  (-b + sqrt(D))/2/a , ( -b - sqrt(D))/2/a  );
				else 
				throw std::exception("D < 0 error!");
			}
			
			int main()
			{
				setlocale(LC_ALL, "Russian");
				int a,b,c;
				std::cin>>a>>b>>c;
				auto [X1,X2] = quadratic_equation(a,b,c);
				std::println("x1={} x2={}",X1, X2);				
				return 0;
			}
			
		\end{cppcode}
		
	\end{block}
	
\end{frame}


\begin{frame}[fragile]\frametitle{Соберём и запустим её}
	
	\begin{block}{CMakeLists.txt}
		
		\begin{cmakecode} 
			cmake_minimum_required(VERSION 4.2)
			project ("quadratic equation")
			set (CMAKE_CXX_STANDARD 23) #Нам нужна современная версия языка для std::println
			add_executable(qe main.cpp)
		\end{cmakecode}
		
		
	\end{block}
	
	\medskip
	
	\begin{Verbatim}[bgcolor=black,tabsize=4,formatcom=\color{white},gobble=2,fontsize=\small, commandchars=\\\{\}, xleftmargin=0em, frame=single,framerule=0pt ,rulecolor=\color{red},framesep=1em,fillcolor=\color{black}]
		
		PS ...> \color{orange}cmake -B build 
		-- Build files have been written to: .../build
		PS ...> \color{orange}cmake --build build
		qe.vcxproj - > .../build/Debug/qe.exe
		\color{orange}1 -5 6  \color{red} <== Вводим коэф. уравнения
		x1=3 x2=2    \color{cyan} <== Все работает!
		
	\end{Verbatim}
	
\end{frame}



\section[.h и .cpp]{Разделение кода на {\ttfamily .h} и {\ttfamily .cpp} файлы }
\frame{\sectionpage}



\begin{frame}[fragile]\frametitle{Основные форматы файлов}
	
	\begin{alertblock}{Заголовочные файлы (с расширением .h, .hpp, .hxx) ---}
		это файлы, которые содержат объявления (декларации) функций, классов, переменных, констант и других сущностей, но не содержат их реализацию.
	\end{alertblock}	
	
	\smallskip
	
	\begin{alertblock}<2->{Файлы кода (с расширением .cpp, .cc, .cxx) ---}
		это файлы, которые содержат реализацию (определения) функций, методов классов и других сущностей, объявленных в заголовочных файлах.
	\end{alertblock}
	
	
	
	\vspace*{-\baselineskip}
	\begin{columns}[]
		\edef\colA{0.5}
		\begin{column}[t]{\colA\linewidth}
			%Колнка1
			\begin{block}{func.h}
				\begin{cppcode}
					void foo();//прототип
					
					class Stankin //описание класса
					{             //и его функций 
						int students_count = 0;
						public:
						void bar();
					}
				\end{cppcode}
			\end{block}
		\end{column}
		\begin{column}[t]{\fpeval{1-\colA}\linewidth}
			%Колнка2
			\begin{block}<2->{func.cpp}
				\begin{cppcode}
					#include<func.h>
					void foo() //реализация foo()
					{
						std::println("fooooo");
					}					
					void Stankin::bar() //реализация 
					{                   //Stankin::bar()
						std::println("baaarrr");
					}
				\end{cppcode}
			\end{block}
		\end{column}
	\end{columns}
	
	
	
\end{frame}


\begin{frame}[fragile]\frametitle{Сделаем новую библиотеку для исходной программы}
	
	\begin{block}{mymath.h}
		\begin{cppcode}
			#include <tuple>
			std::tuple<double,double> quadratic_equation(double, double, double);
		\end{cppcode}
	\end{block}
	
	\begin{block}{mymath.cpp}
		\begin{cppcode}
			#include <cmath>
			#include <exception>
			#include "mymath.h"
			
			std::tuple<double,double> quadratic_equation(double a, double b, double c)
			{
				double D = b*b-4*a*c;
				if (D>0)
				return std::make_tuple(  (-b + sqrt(D))/2/a , ( -b - sqrt(D))/2/a  );
				else 
				throw std::exception("D < 0 error!");
			}
		\end{cppcode}
	\end{block}
	
\end{frame}

\begin{frame}[fragile]\frametitle{Немного подправим main.cpp и CMakeLists.txt}
	
	\begin{block}{main.cpp}
		\begin{cppcode}
			//include <print> <iostream> <clocale>
			#include "mymath.h" // <-Наша библиотека
			//cmath tuple exception   больше не нужны тут
			
			int main()
			{
				setlocale(LC_ALL, "Russian");
				int a,b,c;
				std::cin>>a>>b>>c;
				auto [X1,X2] = quadratic_equation(a,b,c);
				std::println("x1={} x2={}",X1, X2);
				
				return 0;
			}
		\end{cppcode}
	\end{block}
	
	\begin{block}{CMakeLists.txt}
		\begin{cmakecode}
			cmake_minimum_required(VERSION 4.2)
			project ("quadratic equation")
			set (CMAKE_CXX_STANDARD 23)
			add_executable(qe main.cpp mymath.cpp) #Добавляем сюда все cpp
		\end{cmakecode}
	\end{block}
	
\end{frame}

\begin{frame}[fragile]\frametitle{Проверяем результат}
	
	\begin{Verbatim}[bgcolor=black,tabsize=4,formatcom=\color{white},gobble=2,fontsize=\small, commandchars=\\\{\}, xleftmargin=0em, frame=single,framerule=0pt ,rulecolor=\color{red},framesep=1em,fillcolor=\color{black}]
		
		PS ...> \color{orange}cmake -B build
		...
		-- Build files have been written to: .../build
		
		PS ...> \color{orange}cmake --build build/
		...
		qe.vcxproj -> .../build/Debug/qe.exe
		
		\pause{}PS ...> \color{orange} ./build/Debug/qe.exe
		1 -5 6
		x1=3 x2=2 \color{red} <== все работает! 
		\pause{}
		
		PS ...>  \color{orange} tree ./build/qe.dir/Debug /F
		C:/.../BUILD/QE.DIR/DEBUG
		│   main.obj  				\color{cyan} <== наши obj
		│   mymath.obj              \color{cyan} <== наши obj
		│   ...
		...
		
	\end{Verbatim}
	
\end{frame}

\begin{frame}  \frametitle{Схема включений}
	\centering
	
	\begin{tikzpicture}[
		file/.style={draw, font=\ttfamily},
		lib/.style={draw, font=\ttfamily\itshape\color{ForestGreen},}
		]
		\draw node[file] (maincpp) {main.cpp};
		
		
		\draw (maincpp) ++ (0:3) node[file] (mymathcpp) {mymath.cpp}; 
		
		\draw (mymathcpp) ++(90:2) node[file] (mymathh) {mymath.h};
		
		\draw (mymathh) +(90:2) node[lib] (typle) {<typle>};
		
		\draw (maincpp) +(75:3) node[lib] (print) {<print>}
		+(120:3) node[lib] (clocale){<clocale>}
		+(150:3) node[lib] (exception) {<exception>};
		
		\draw[-Latex] 	(mymathh) -- (maincpp)  ;
		\draw[-Latex]	(mymathh) -- (mymathcpp);
		\draw[-Latex]	(typle) -- (mymathh);
		\draw[-Latex]   (print) -- (maincpp) ;
		\draw[-Latex]   (clocale) -- (maincpp) ;
		\draw[-Latex]   (exception) -- (maincpp) ; 	
		
	\end{tikzpicture}
\end{frame}


\begin{frame}[fragile]\frametitle{Усложним проект}
	
	\begin{block}{mymath.h}
		\begin{cppcode}
			#include <tuple>
			std::tuple<double,double> quadratic_equation(double, double, double);
			
			@\rememb{a}@bool isEven(int n) 
			{
				return n % 2 == 0;
			}                    @\rememb{b}@    
		\end{cppcode}
	\end{block}
	
	\begin{tikzpicture}[remember picture, overlay]
		\node[draw, red, fit=(a)(b)] {};
	\end{tikzpicture}
	
	\begin{block}{main.cpp}
		\begin{cppcode}
			...
			int main()
			{
				...
				
				@\rememb{a}@std::println( "x1 is {}", isEven(X1)?"even":"odd" );@\rememb{b}@    
				
				return 0;
			}
		\end{cppcode}
	\end{block}
	
	\begin{tikzpicture}[remember picture, overlay]
		\node[draw, red, fit=(a)(b)] {};
	\end{tikzpicture}
	
\end{frame}

\begin{frame}[fragile]\frametitle{Собираем}
	
	\begin{Verbatim}[bgcolor=black,tabsize=4,formatcom=\color{white},gobble=2,fontsize=\small, commandchars=\\\{\}, xleftmargin=0em, frame=single,framerule=0pt ,rulecolor=\color{red},framesep=1em,fillcolor=\color{black}]
		
		PS ..> \color{orange}cmake --build build/
		Версия MSBuild 18.3.0-release-26070-10+3972042b7 для .NET Framework
		
		Проверка источников на наличие зависимостей модуля…
		mymath.cpp
		main.cpp
		Компиляция…
		main.cpp
		mymath.cpp
		...
		\color{red}mymath.obj : error LNK2005: "bool __cdecl isEven(int)" (?isEven@@YA_NH@Z) 
		\color{red}уже определен в main.obj
		\color{red}fatal error LNK1169: обнаружен многократно определенный
		\color{red}символ - один или более 
		
		
	\end{Verbatim}
	
\end{frame}

\begin{frame}[fragile]\frametitle{Решение проблемы --- inline}
	
	\begin{block}{mymath.h}
		\begin{cppcode}
			#include <tuple>
			std::tuple<double,double> quadratic_equation(double, double, double);
			
			@\rememb{a}@inline@\rememb{b}@  bool isEven(int n) 
			{
				return n % 2 == 0;
			}                     
		\end{cppcode}
	\end{block}
	
	\begin{tikzpicture}[remember picture, overlay]
		\node[draw, red, fit=(a)(b)] {};
	\end{tikzpicture}
	
	\begin{Verbatim}[bgcolor=black,tabsize=4,formatcom=\color{white},gobble=2,fontsize=\small, commandchars=\\\{\}, xleftmargin=0em, frame=single,framerule=0pt ,rulecolor=\color{red},framesep=1em,fillcolor=\color{black}]
		
		PS ...> \color{orange}cmake --build build/
		...
		qe.vcxproj -> C:.../build/Debug/qe.exe
		
		PS ...> \color{orange}./build/Debug/qe.exe
		1 -9 20
		x1=5 x2=4
		x1 is not odd \color{red}<==
		
	\end{Verbatim}
	
	
\end{frame}

\begin{frame}[fragile]\frametitle{Встроенные функции (inline)}
	
	\begin{alertblock} {Inline (подстановка)  ---}
		это указание компилятору, что код функции нужно вставить непосредственно в место вызова, вместо генерации инструкции вызова (call).
	\end{alertblock}
	
	\textbf{Основные цели}
	
	\begin{itemize}
		\item Оптимизация скорости — убираются накладные расходы на вызов функции (создание стекового кадра, передача аргументов, возврат)
		
		\item Использование в заголовочных файлах — позволяет определять функции в .h файлах без ошибок множественного определения
	\end{itemize}
	
	\begin{columns}[onlytextwidth]
		\edef\colA{0.5}
		\begin{column}{\colA\linewidth}
			%Колнка1
			\begin{cppcode}
				// Без inline
				int add(int a, int b) {
					return a + b;
				}
				
				int main() {
					int x = add(5, 3);
					// Здесь будет инструкция CALL  
					return 0;
				}
			\end{cppcode}
		\end{column}
		\begin{column}{\fpeval{1-\colA}\linewidth}
			%Колнка2
			\begin{cppcode}
				// C inline
				inline int add(int a, int b) {
					return a + b;
				}
				
				int main() {
					int x = add(5, 3);
					// Компилятор заменит на: int x = 5 + 3;  
					return 0;
				}
			\end{cppcode}
		\end{column}
	\end{columns}
	
	
	
\end{frame}


\begin{frame}[fragile]\frametitle{Добавим константу в библиотеку}
	
	\begin{block}{mymath.h}
		\begin{cppcode}
			#include <tuple>
			std::tuple<double,double> quadratic_equation(double, double, double);
			
			inline bool isEven(int n) {
				return n % 2 == 0;
			}
			
			//не будет работать!!
			double half_PI = 1.5707963268; 
		\end{cppcode}
	\end{block}
	
	\pause
	
	\textbf{Варианты:}
	
	\begin{cppcode}
		//1. будет
		const double half_PI = 1.5707963268; 
		//2. будет с++11
		constexpr double half_PI = 1.5707963268; 
		//3. будет даже так!  с++17
		inline double half_PI = 1.5707963268; 
		//4. будет, но не стоит
		static double half_PI = 1.5707963268; 
	\end{cppcode}
	
	
\end{frame}


\begin{frame} [fragile]  \frametitle{Добавим новую библиотеку}
	\begin{columns}[onlytextwidth]
		\edef\colA{0.75}
		\begin{column}{\colA\linewidth}
			%Колнка1
			\begin{tikzpicture}[
				file/.style={draw, font=\ttfamily},
				lib/.style={draw, font=\ttfamily\itshape\color{ForestGreen},}
				]
				\draw node[file] (maincpp) {main.cpp};
				
				
				\draw (maincpp) ++ (0:3) node[file] (mymathcpp) {mymath.cpp}; 
				
				\draw (mymathcpp) ++(90:2) node[file] (mymathh) {mymath.h};
				
				\draw (mymathh) +(90:2) node[lib] (typle) {<typle>};
				
				\draw (maincpp) +(110:5) node[lib] (print) {<print>}
				+(130:5) node[lib] (clocale){<clocale>}
				+(150:5) node[lib] (exception) {<exception>};
				
				\draw (maincpp) +(80:4) node[file,red] (pointh) {point.h};
				
				\draw[-Latex] 	(mymathh) -- (maincpp)  ;
				\draw[-Latex]	(mymathh) -- (mymathcpp);
				\draw[-Latex]	(typle) -- (mymathh);
				\draw[-Latex]   (print) -- (maincpp) ;
				\draw[-Latex]   (clocale) -- (maincpp) ;
				\draw[-Latex]   (exception) -- (maincpp) ; 	
				\draw[-Latex, red]   (pointh) -- (maincpp) ; 	
				\draw[-Latex, red]   (pointh) -- (mymathh) ; 
				
			\end{tikzpicture}
		\end{column}
		\begin{column}{\fpeval{1-\colA}\linewidth}
			%Колнка2
			\begin{block}{Point.h}
				\begin{cppcode}
					struct Point
					{
						double x;
						double y;
					};
				\end{cppcode}
			\end{block}
		\end{column}
	\end{columns}
	
	
\end{frame}


\begin{frame}[fragile]\frametitle{Собираем с новой библиотекой}
	
	\begin{Verbatim}[bgcolor=black,tabsize=4,formatcom=\color{white},gobble=2,fontsize=\small, commandchars=\\\{\}, xleftmargin=0em, frame=single,framerule=0pt ,rulecolor=\color{red},framesep=1em,fillcolor=\color{black}]
		
		PS ...> \color{orange} cmake --build build/
		Версия MSBuild 18.3.0-release-26070-10+3972042b7 для .NET Framework
		
		Проверка источников на наличие зависимостей модуля…
		mymath.cpp
		main.cpp
		Компиляция…
		mymath.cpp
		main.cpp
		C:/.../point.h(1,8): 
		\color{red}error C2011: Point: переопределение типа "struct" 
		\color{red}[C:.../build/qe.vcxproj]
		\color{red}(компиляция исходного файла "../main.cpp")
		\color{red}C:/.../point.h(1,8):
		\color{red}см. объявление "Point"
		
	\end{Verbatim}
	
\end{frame}

\begin{frame}[fragile]\frametitle{Стражи заголовков}
	
	\begin{alertblock}{Стражи заголовков (Include Guards) --- }
		конструкция препроцессора, защищающая от повторного включения одного и того же заголовочного файла. Предотвращают множественные определения функций и классов, которые вызывают ошибки компиляции.
	\end{alertblock}
	
	\pause\textbf{Два способа реализации:}
	
	\begin{columns}[]
		\edef\colA{0.45}
		\begin{column}[t]{\colA\linewidth}
			%Колнка1
			
			Через директиву \verb|#pragma once| 
			
			\begin{block}{Point.h}
				\begin{cppcode}
					#pragma once //просто, 
					//но не стандартно
					struct Point
					{
						double x;
						double y;
					};
				\end{cppcode} 
				
			\end{block}  
		\end{column} \pause
		\begin{column}[t]{\fpeval{1-\colA}\linewidth}
			%Колнка2	
			
			Через конструкцию \verb|#ifndef ... #endif| 
			
			\begin{block}{Point.h}
				\begin{cppcode}
					#ifndef POINT_H //сложнее, но всегда работает
					#define POINT_H //нужно придумать такое имя,
					//которое никогда не повторится
					struct Point
					{
						double x;
						double y;
					};
					
					#endif //POINT_H
				\end{cppcode}
			\end{block}
		\end{column}
	\end{columns}
	
\end{frame}


	
\end{document}
